\scanpage{42}
\noindent\hfill atkin \mspage{2}

\noindent with Dr. R.F.Churchhouse,both a number theorist and a veteran programmer,and Mr. \handwritten{E}.B.Fossey and Miss Barbara Stokoe have constantly assisted me with the Atlas Hartran(Fortran) system.

It is intended that this account shall be intelligible to the reader with no knowledge of modular theory\deleted{;}\margin{\handwritten{\figure{\tikz[x=1ex,y=1ex,baseline=0pt,line cap=round]{\draw (-.15,-.65)--(1.05,1.2);\draw (.52,1.0).. controls (-.28,.9) and (-.25,.35)..(.2,.15).. controls (.55,-.03) and (.5,-.35)..(.18,-.5);}}\ \inserted{,}\ (comma)}}provided that he is prepared to accept that the coefficient \uline{$c(n)$} which we mainly consider is mathematically significant. However at some points we give brief additional explanations for the reader who has such knowledge.

\noindent\margin{\handwritten{/}}2. \uline{Ramanujan} \deleted{C}\uline{ongruences} \uline{for} \uline{primes} $q\handwritten{\leq}11$. The general theory of elliptic and modular functions had been in the process of development for about a century before RAMANUJAN(1919)conjectured his remarkable congruences for the partition function \uline{$p(n)$},defined combinatorially as the number of distinct ways of expressing \uline{$n$} as the sum of integral parts,where the order of the parts is immaterial. His conjecture was
\[
\handwritten{\text{\uline{if}}\quad 24n-1\equiv0\;(\text{\uline{mod}}\;5^{a}7^{b}11^{c})\quad
\text{\uline{then}}\quad p(n)\equiv0\;(\text{\uline{mod}}\;5^{a}7^{b}11^{c})}
\tag*{\handwritten{$\ldots$ (2.1)}}
\]
The curious appearance of $24n-1$ in \handwritten{(2.1)} is due to the fact that, writing
\[
\begin{aligned}
&\handwritten{x=e^{2\pi i\tau}\qquad (\,|x|<1,\ \mathrm{Im}\,\tau>0),}\\
&\handwritten{f(x)=(1-x)(1-x^{2})(1-x^{3})\ldots,}\\
\text{\deleted{we have}}\quad &\handwritten{\eta(\tau)=e^{\pi i\tau/12}f(x),}\\
\text{\inserted{\handwritten{we have}}}\quad &\handwritten{\sum_{n=0}^{\infty}p(n)\,x^{n}=1/f(x),}
\end{aligned}
\tag*{\handwritten{(2.2)}}
\]
and $\handwritten{\eta(\tau)}$ is a modular form of half-integral dimension and complicated multiplier system. Thus the congruence properties are naturally associated not with \uline{$f(x)$} whose reciprocal enumerates \uline{$p(n)$},but with \uline{$x^{1/24}f(x)$} ; that is,they are essentially modular and not combinatorial.$^{*}$ Ramanujan himself was of course well aware of this,

\noindent\handwritten{\rule{\linewidth}{0.4pt}}

\noindent\deleted{as his proofs}

\noindent *But see DYSON(19\handwritten{44}) and ATKIN and SWINNERTON-DYER(1954)for an exception.

\noindent\handwritten{\rule{\linewidth}{0.4pt}}

\noindent as his proofs of some particular cases show.

However this conjecture \handwritten{(2.1)},originally found from the

% archive sheet number: 41
